\chapter{Differential Equations}

Differential equations are mathematical equations that relate a function to its derivatives. They are fundamental in describing various physical phenomena, such as motion, heat transfer, and population dynamics. In this chapter, we will explore the basics of differential equations, including their classification, methods of solution, and applications. There are two main types of differential equations: 
\begin{itemize}
    \item \ul{Ordinary Differential Equations (ODEs)}: These involve functions of a single variable and their derivatives.

    \item \ul{Partial Differential Equations (PDEs)}: These involve functions of multiple variables and their partial derivatives.
\end{itemize}

\begin{definicion}[Order of a Differential Equation]
The order of a differential equation is the highest order of derivative that appears in the equation.
\end{definicion}

Some common examples of ODEs include:
\begin{itemize}
    \item The simple harmonic oscillator. According to Hooke's Law, the restoring force $F_R$ exerted by a spring is proportional to the displacement $x$ of the mass from its equilibrium position:
    \[
        F_R = -kx
    \]
    where $k$ is the spring constant. According to Newton's second law, the force is also equal to the mass $m$ times the acceleration (the second derivative of displacement with respect to time):
    \[
        F_R = m \frac{d^2x}{dt^2}
    \]
    Combining these two equations gives us the second-order ordinary differential equation:
    \[
        \frac{d^2x}{dt^2} + \frac{k}{m}x = 0
    \]

    \item The Radiative Decay Law. This law ensures that the rate of decay of a radioactive substance is inversely proportional to the amount of substance present. Mathematically, this can be expressed as:
    \[
        \frac{dN}{dt} = -\lambda N
    \]
    where $N(t)$ is the number of radioactive atoms at time $t$, and $\lambda$ is the decay constant.

    \item Free fall under gravity. The motion of an object in free fall can be described by the second-order ordinary differential equation:
    \[
        \frac{d^2y}{dt^2} = -g
    \]
    where $y(t)$ is the height of the object at time $t$, and $g$ is the acceleration due to gravity.

    \item When working with capacitors and inductors in electrical circuits, we often encounter first-order ODEs.
\end{itemize}

It should be noted that a ODE does not usually have a unique solution, as it may contain arbitrary constants that can take on any value. To determine a unique solution, we need to specify initial conditions or boundary conditions.

\section{ODE solution methods}

To solve ordinary differential equations, we can use various methods depending on the type and complexity of the equation.

\subsection{Primitive finding}

The simplest method to solve an ODE is to find its primitive. This works for first-order ODEs of the form:
\[
    \frac{dy}{dx} = f(x)
\]
In this case, we can integrate both sides with respect to $x$ to find the solution:
\[
    y(x) = \int f(x) \ dx + C
\]
where $C$ is the constant of integration.

\subsection{Separation of variables}

For first-order ODEs that can be expressed in the form:
\[
    \frac{dy}{dx} = g(x)h(y)
\]
we can use the method of separation of variables. This involves rearranging the equation to separate the variables:
\[
    \int \frac{1}{h(y)} \ dy = \int g(x) \ dx
\]

This allows us to integrate both sides independently to find the solution. The solution is usually expressed in implicit form, and we may need to solve for $y$ explicitly if possible.

\section{Types of ODE}

Once we have a general understanding of how to solve ODEs, we can classify them into different types based on their structure and properties. We will then explore specific methods for solving each type of ODE.

\subsection{Linear ODE of first order}

A first-order linear ODE has the form:
\[
    \frac{dy}{dx} + P(x)y = Q(x)
\]

In order to solve this type of equation, a distinction is made between the homogeneous case, where $Q(x) = 0$, and the non-homogeneous case, where $Q(x) \neq 0$. 

\subsubsection{Homogeneous case}

In the homogeneous case, the equation simplifies to:
\[
    \frac{dy}{dx} + P(x)y = 0
\]
As we can see, this is a separable equation. Solving it gives us the general solution:
\[
    y(x) = Ce^{\displaystyle -\int P(x) \ dx}
\]

\subsubsection{Non-homogeneous case}
In the non-homogeneous case, there are two approaches to find the solution:
\begin{itemize}
    \item If we can find a particular solution $y_p(x)$ to the non-homogeneous equation, then the general solution can be expressed as:
    \[
        y(x) = y_h(x) + y_p(x)
    \]
    where $y_h(x)$ is the general solution to the corresponding homogeneous equation. Finding $y_p(x)$ may not be easy, but there are tables that give hints on how can it be.

    \item Alternatively, once $y_h(x)$ is known, the constant $C$ in the homogeneous solution can be replaced with a function $K(x)$, leading to the method of variation of constant. This involves substituting
    \[
        y(x) = K(x)e^{\displaystyle -\int P(x) \ dx}
    \]
    into the original non-homogeneous equation to derive an equation for $K(x)$, which can then be solved to find its value. Once $K(x)$ is determined, the general solution to the non-homogeneous equation is also knwon.
\end{itemize}

\subsection{Linear ODE with constant coefficients}

A linear ODE with constant coefficients has the form:
\[
    \frac{d^n y}{dx^n} + a_{n-1} \frac{d^{n-1} y}{dx^{n-1}} + \cdots + a_1 \frac{dy}{dx} + a_0 y = g(x)
\]
where $a_0, a_1, \ldots, a_{n-1}$ are constants and $g(x)$ is a function of $x$.

\subsubsection{Homogeneous case}

In the homogeneous case, where $g(x) = 0$, the solutions form a vector space of dimension $n$. To find the general solution, we can use the characteristic equation:
\[
    \lambda^n + a_{n-1} \lambda^{n-1} + \cdots + a_1 \lambda + a_0 = 0
\]
The roots of this characteristic equation determine the form of the general solution.
\begin{itemize}
    \item For every distinct (real or complex) root $\lambda_i$, there is a corresponding solution of the form:
    \[
        y_i(x) = e^{\lambda_i x}
    \]
    \item For every repeated root $\lambda$ of multiplicity $m$, there are $m$ linearly independent solutions of the form:
    \[
        y_i(x) = e^{\lambda x}, \quad y_{i+1}(x) = xe^{\lambda x}, \quad \ldots, \quad y_{i+m-1}(x) = x^{m-1}e^{\lambda x}
    \]
    \item For every complex solution $z(t)$ of the ODE, two real solutions are known:
    \[
        y_1(t) = \Re(z(t)), \quad y_2(t) = \Im(z(t))
    \]
\end{itemize}

That way, we find $n$ different real solutions. In order to check that they are linearly independent, we can compute the Wronskian determinant.
\[
    W(y_1, y_2, \ldots, y_n)(t) = \begin{vmatrix}
    y_1(t) & y_2(t) & \cdots & y_n(t) \\
    y_1'(t) & y_2'(t) & \cdots & y_n'(t) \\
    \vdots & \vdots & \ddots & \vdots \\
    y_1^{(n-1)}(t) & y_2^{(n-1)}(t) & \cdots & y_n^{(n-1)}(t)
    \end{vmatrix}\qquad t\in I
\]
If the Wronskian is non-zero for some $t$ in the interval $I$, then the solutions are linearly independent. Therefore, the general solution to the homogeneous equation can be expressed as a linear combination of these independent solutions.

\subsubsection{Non-homogeneous case}

In the non-homogeneous case, where $g(x) \neq 0$, variation of constant could be used to find the general solution. Alternatively, if $g(x)$ is of a specific form (e.g., polynomial, exponential, sine, or cosine), we can use the method of undetermined coefficients to find a particular solution $y_p(x)$. Once we have $y_p(x)$, the general solution to the non-homogeneous equation can be expressed as:
\[
    y(x) = y_h(x) + y_p(x)
\]